\documentclass[oneside, 12pt]{book}

\usepackage[ukrainian]{babel}
\usepackage[T2A]{fontenc}
\usepackage[utf8]{inputenc}
\usepackage{indentfirst}
\usepackage{titling}
\usepackage[pagestyles]{titlesec}

\renewcommand{\bfseries}{}
\setcounter{tocdepth}{1}

\titleformat{\section}[display]{}{}{0pt}{\LARGE}
\newpagestyle{mystyle}
{\sethead{}{}{\thepage}}
\pagestyle{mystyle}

\newenvironment{prat}[1]
{\begin{center}
\begin{tabular}{p{0.9\textwidth}}
\hline\hline
\multicolumn{1}{c}{Запрошення до практики}\\
\hline
Частота: #1.\\
\hline
}{\\
\hline\hline
\end{tabular} 
\end{center}}

\title{Про досконалість Буття}
\author{Артур Пауліну}
\date{}

\begin{document}

\maketitle

\chapter*{Правове застереження}
\thispagestyle{empty}

Скориставшись будь-якою інформацією, поданою в цій книзі, читач цілковито відповідає сам за себе.

\chapter*{Подяки}
\thispagestyle{empty}

Я глибоко вдячний своїм батькам за підтримку навіть у найважчі миті; усім просвітленим істотам, які вже існували, існують і існуватимуть, і моїй супутниці Ларіссі, яка є справжньою окрасою.

\tableofcontents
\thispagestyle{empty}

\chapter*{Передмова}\addcontentsline{toc}{chapter}{Передмова}

Рекомендую це читання тим, хто вже відчув тугу духовних злиднів і зумів утвердити в собі бажання пережити Життя в його найповнішій повноті. Отже, це не книга, яку читають наспіх. Кожне речення --- частина продуманої аргументації, покликаної запросити читача до самозаглиблення й заохотити його до глибокої трансформації в напрямку майстерності над собою.

Мова буде проста й предметна, і часто я говоритиму від першої особи. Роблю так, щоб твір був ще чеснішим і щирішим, адже тут присутні також мої почуття вдячності й відповідальності за людство здоровіше, заможніше й ближче до усвідомлення своєї єдності.

\chapter*{Вступ}\addcontentsline{toc}{chapter}{Вступ}

Як нам жити добре, коли є так багато причин скаржитися? Джерел незадоволення практично не злічити. Нам не подобається соціальний, політичний та економічний стан, у якому ми перебуваємо, і ми зазвичай звинувачуємо в цьому наших так званих лідерів або навіть брак вибору на виборах. Але вже час нам шукати глибшого погляду й розуміти силу, яку ми маємо як окремі особи.

Це факт, що політики цієї доби людства приходять до влади за допомогою капіталу і з наміром йому служити. Через жадібність великих корпорацій при владі ми приречені на дедалі поляризованіше зростання на користь небагатьох і на руйнування екосистем цієї планети.

Нам треба зрозуміти корінь проблеми. Зрештою, звідки береться влада цих великих злонамірених корпорацій? Відповідь на це питання --- чистий нектар: від нас! Від нас із нашими вадами, що спонсорують незнання, і воно набуває світового масштабу.

Безвідповідальність щодо походження нашої їжі, сидячий спосіб життя, пошук задоволення у вживанні одурманювальних речовин, поширення відволікання, примусове споживацтво, вписане в культуру показності, поширення насильства, пісні й серіали, що проповідують прив'язаність, гнів і заздрість, переоцінювання екстравертності, пошук поверховних стосунків, нарцисизм, дисбаланс між працею й родинною близькістю. Усе це --- елементи, що творять дедалі хаотичніше й нелюдяніше суспільство. І щоб звільнитися від усього цього, нам треба пливти проти течії зручності й виконати найважче завдання: знайти в собі джерело туги, що штовхає нас до цих джерел безладу.

Необхідна індивідуальна праця --- це очищення. Кожна людина має звільнитися від взірців, що завдають страждання їй самій і ближньому. Однак якщо ми ведемо розсіяне життя, ми стаємо нездатними бачити такі джерела страждання в собі й просто живемо у владі своєї отруйності, випромінюючи її на всіх істот.

Шлях, який я бачу для нашого зцілення, --- це самопізнання. Це і є спонукання до появи цієї книги. Запропоновані в цьому творі студії не схожі на ті, що ведуть у промисловій школі. Ми не оглядатимемо теми, щоб виробити холодний і відсторонений від нас погляд, ніби маємо справу з чимось цілком чужим, а навпаки --- якнайінтимніший.

Труднощі й легкість навчання духовності полягають у таких фактах: \textit{i}) мудрість, проповідувана й не практикована, лицемірна й недієва; \textit{ii}) мудрість, проповідувана й практикована, надихає й приносить користь. Тож я приходжу з правом говорити про практики, на яких я тримався чи тримаюся, і з обов'язком не перевищувати їх ні в чому.

Структура викладена в чотирьох розділах, які розглядають теми, пов'язані з ділянками типового усвідомленого переживання, у зростаючому порядку дрібності: від найгрубшого до найтоншого.

\begin{itemize}
\item Тіло --- носій, посередник досвіду
\item Серце --- метафора, що охоплює почуття
\item Розум --- абстрактний світ думок
\item Трансцендентне --- по той бік мови, Істина
\end{itemize}

Мета цього поділу --- не цілковите відокремлення Реальності на незалежні частини. Постава тіла впливає на поведінку Розуму, а думки можуть запускати певні емоції. Або ж емоції живлять думки, що можуть спричиняти зміни в стані Тіла. А Трансцендентне своєю чергою пронизує все існування.

Ця книга містить численні вправи, і якщо читач не має готовності постійно їх практикувати, немає потреби продовжувати читання, адже духовне самовдосконалення --- це процес безперервного тренування.

Можна буде помітити, що багато нагод відповідають більш ніж одній вправі. У таких випадках читач сам обирає вправу, яка найкраще його скеровує.

\chapter*{Про досконалість Тіла}\addcontentsline{toc}{chapter}{Про досконалість Тіла}

Ми можемо з'єднатися з Тілом і навчитися відчувати та поважати його сигнали. Піддаватися невпорядкованому режимові й сподіватися, що Тіло його витримає, --- це акт насильства проти самої природи, що дає довгу низку розладів. Не випадково знеболювальні є частиною нашої культури.

Для Досконалості Тіла я виокремив чотири аспекти в порядку нагальності: дихання, гідратація, харчування й постава.

За останні два століття наша медицина значно просунулася в розумінні основної механіки тканин Тіла. Надзвичайно важливо звертатися по фахову допомогу, коли це потрібно. Але щодо профілактичної медицини, то її ми здатні втілити у своє життя.

\section*{Дихання}\addcontentsline{toc}{section}{Дихання}

Дихання не лише найнагальніший процес, а й безцінний ключ до самопізнання. Через його величезну вагу до нього йтиметься ще не раз у цій книзі.

Немовлятами ми здебільшого дихаємо цілком повно: грудний, черевний, повільний і широкий. Але навіть у новонароджених дихання стає коротшим і поверховішим у миті неспокою. Дорослішаючи, ми дозволяємо дедалі більшій кількості нездорових елементів проникати в наше життя, аж поки, дорослими, не маємо жодної миті спокою, бо цілком зумовлені вкрай убогим диханням. Використання нижньої частини легень збільшує дихальну ємність, а повний видих усуває залишковий вуглекислий газ.

\begin{prat}{якнайчастіше}
-- Не втручаючись, просто спостерігайте, як відбувається дихання, і усвідомте його ритм і глибину.
\end{prat}

\section*{Гідратація}\addcontentsline{toc}{section}{Гідратація}

Важко назвати хімічну реакцію Тіла, яка не відбувалася б у водному середовищі, тож якісна гідратація неодмінна.

\begin{prat}{якнайчастіше}
-- Усвідомте частоту спраги;

-- Спостерігайте зв'язок між кількістю випитої води, станом фізичної втоми, кількістю виділеної сечі та її кольором;

-- Усвідомте відчуття дискомфорту під час сечовипускання.
\end{prat}

\section*{Харчування}\addcontentsline{toc}{section}{Харчування}

Акт харчування втрачає вагу щодо «важливих справ» дня аж до того, що стає відволіканням. Дедалі більше людей пристає до політики з'їсти будь-що, аби вгамувати голод, або в пошуку приємного відчуття надмірно повного шлунка. І щоб не гаяти часу за їжею, вони виконують якусь іншу справу, лишаючи акт їжі на другому плані.

Пропозиція тут --- не радити читачеві змінити дієту, а скерувати його відчути, як такі процеси відбуваються в Тілі.

Є страви, що спричиняють піки задоволення, за якими йдуть спади настрою, стимулюючи цикли психологічної залежності; що перевантажують шлунок і кишківник, викликаючи небажане нездужання й отруєння; що привчають тіло до неприродної рівноваги, творячи хімічні залежності; що запускають запальні процеси в різних тканинах тощо. Тоді як інші дають легкість настрою, легко засвоюються і не викликають ні залежностей, ні незручностей.

\begin{prat}{перед, під час і після кожного приймання їжі}
-- Перевірте, чи бажання їсти справді походить від голоду;

-- Смакуйте кожен аспект акту їжі;

-- Усвідомте думки, що з'являються під час трапези;

-- Спостерігайте відчуття в шлунку й кишківнику після трапези.
\end{prat}

\section*{Постава}\addcontentsline{toc}{section}{Постава}

Бути відповідальним за напрямок розбудови Тіла --- це довгострокова інвестиція надзвичайної ваги для здорового старіння. Через те що Тіло постійно будує дедалі більше себе, ніколи не пізно почати.

Основна властивість Тіла, здатного дарувати велике самопочуття, --- це гнучка, бадьора мускулатура, здатна розслаблятися в неактивності, особливо коли йдеться про групи м'язів, пов'язані з підтриманням хребта.

Є чимало фахівців у галузі здоров'я, компетентних скерувати читача в цій справі. Серед них фізіотерапевти, ортопеди, інструктори пілатесу тощо. Зокрема, фахівець, якого я рекомендую найбільше, --- це добрий учитель йоги. Але ось досить потужна вправа, здатна спричинити глибокі зміни.

\begin{prat}{якнайчастіше}
-- Випряміть хребет і розкрийте грудну клітку;

-- Пройдіть тілом фокусом уваги й розслабте кожен м'яз, особливо м'язи плечей і обличчя.
\end{prat}

\chapter*{Про досконалість Серця}\addcontentsline{toc}{chapter}{Про досконалість Серця}

Ми живемо в добу, коли емоційний біль не шанують належно, а будь-яка ознака психологічного дискомфорту стає приводом для галасу. Люди почуваються винними, переконані, що чинять помилку, сумуючи.

Але будьмо розважливі. Якщо ми рішуче не бажаємо ставати на шлях насильства, нам треба навчитися жити з тим, що відчуваємо.

У цьому розділі ми приділимо увагу своїм почуттям і перетворимо свої емоції на джерела зцілення й орієнтиру.

\section*{Емоційне дозрівання}\addcontentsline{toc}{section}{Емоційне дозрівання}

Передусім важливо зрозуміти, чим є емоції на практиці. Коли ми тривожимося, можемо відчувати хвилювання в ділянці сонячного сплетення. Коли нам страшно, можемо відчувати холод в основі шлунка. Коли нам дуже сумно, можемо відчувати дискомфорт у горлі. Для практичних цілей вважатимемо емоціями дотикові прояви почуттів у тілі.

Два головні способи затримати емоційне дозрівання --- це спроба вести знеболене життя і плекання жертовної постави. У першому випадку можна побачити занедбаність, у другому --- упертість.

Той, хто занедбує, схильний повторювати ті самі помилки й накопичувати ще більше прихованого страждання. Упертий примножує власний біль, переживаючи ті самі історії знову й знову. Річ у тім, що нам потрібна мужність відчути те, що життя нам подає.

\begin{prat}{двічі на тиждень}
-- Сядьте з прямим хребтом на сідничних горбах і покладіть руки на ноги;

-- Заплющте очі, дихайте усвідомлено й знайдіть емоцію в тілі;

-- Зберігаючи усвідомлене дихання, дайте повну увагу емоційному дискомфорту;

-- Просто відчувайте біль і дозвольте його причині піти.
\end{prat}

\section*{Любов до себе}\addcontentsline{toc}{section}{Любов до себе}

Одного разу людина розповіла мені таку ситуацію: вона телефонувала всім друзям зі свого списку контактів, щоб спитати, як вони, і побажати їм щастя, але почувалася сумною, бо ніхто не робив того самого для неї. Тоді я попросив її якийсь час бажати щастя своїм друзям у тиші, не телефонуючи їм. За кілька днів вона прийшла сказати мені, що не може цього зробити, бо відчуває безмежну самотність. Це чудовий приклад для ілюстрації цієї теми.

Я вже знав, що вона не зможе виконати завдання, яке я їй дав, але дав усе одно, бо їй також треба було дізнатися, що вона не зможе. Вона телефонувала друзям у пошуку визнання, намагаючись заповнити власну потребу. Вона навіть не могла щиро побажати їм щастя.

Що ж, як нам жити любов'ю до себе? Якщо бути належно щирими, ми побачимо, що схильні до різних глибоко вкорінених переконань. Людина з наведеного випадку вірила, що потребує зовнішнього підтвердження, щоб бути щасливою, --- надзвичайно поширене переконання, яке робить приклад ще кращим.

Кожне руйнівне переконання, від якого ми звільняємося, дає більше простору, щоб любов до себе розквітла в нашому житті. Я визнаю, що зіткнутися з власними переконаннями --- справа нелегка. Адже ми десятиліттями їх плекали.

\begin{prat}{раз на тиждень}
-- Випишіть усе, що ви здатні сказати про себе, короткими твердженнями, насамперед те, що стосується необхідного для щастя;

-- Рішуче оспорте кожен пункт, навіть найтривіальніший, і позначте ті, що впадуть перед вашою найпотаємнішою правдою;

-- Протягом тижня усвідомлюйте кожен прояв, пов'язаний із позначеним у списку переконанням.
\end{prat}

\section*{Співчуття}\addcontentsline{toc}{section}{Співчуття}

Одним із моїх перших потрясінь після початку роботи над самотрансформацією було усвідомлення, скільки страждання я завдав навколо себе вчинками, мотивованими власними тривогами. Тоді мене наче блискавкою осягнуло розуміння, що це типовий людський взірець: ми завдаємо страждання, бо страждаємо.

Тож поміркуймо про людей і спробуймо їх зрозуміти. Від найближчих, із якими в нас можуть бути ворожнечі, до найдальших, як-от корумповані політики, --- усі вони схильні до кривд життя. Людство, занурене в незнання, буде великим океаном страждання.

Моя пропозиція --- обрати вихід із багна страждання, зі світу ворожнеч. Свідомо оберімо співчуття, насамперед до тих, хто найбільше шкодить, бо саме вони найбільше страждають. Досліджуймо явище страждання в собі й в інших.

\begin{prat}{раз на день}
-- Подумайте про людину, яка завдає чи завдала вам туги (це можете бути ви самі);

-- Визнайте причини страждання цієї людини, не судячи їх. Якщо можливо, дружня розмова --- найпотужніший спосіб;

-- Побажайте, щоб ця людина змогла звільнитися від страждання і досягла спокою.
\end{prat}

\chapter*{Про досконалість Розуму}\addcontentsline{toc}{chapter}{Про досконалість Розуму}

Розум --- потужний інструмент людського витвору, але величезна більшість із нас досі збита з ніг, бо не вміє орієнтуватися в його просторі. Брак свободи в опануванні Розуму призводить до того, що нас несвідомо несуть руйнівні думки, і життя обмежується фрустраціями й тривогами. Ми не матимемо умов справді процвітати, доки настанови щодо поводження з Розумом не стануть широко відомими в знанні людства.

Як досягти розумового звільнення і згасити цикл незадоволення? Як перестати бути лицемірними, бажаючи миру в світі? Це питання, що ведуть нас до Досконалості Розуму.

\section*{Об'єкти Розуму}\addcontentsline{toc}{section}{Об'єкти Розуму}

Об'єкт Розуму може належати до одного з трьох типів: пам'ять, уява або судження.

Пам'ять і уява мають подібну природу. Зазвичай це образи, запахи, звуки або послідовності подій. Часто ми сумніваємося, чи щось справді сталося, чи це уява.

Судження відбуваються через розумовий голос. Читач може спробувати заплющити очі й подумати фразу «Планета Земля блакитна». Це наче ми чуємо голос, який може навіть викликати об'єкт пам'яті (напр.: фото планети Земля) або уяви (напр.: блакитна куля, що ширяє в космосі).

Перешкода для переживання Досконалості Розуму --- це те, що ми почуваємося безсилими перед потоком думок. Звичайні приклади проявів цієї зумовленості --- постійна втеча в добрі спогади, постійна втеча в уявні насолоди, бажання змінити минуле, одержимість контролем над майбутнім, самобичування й нездатність прощати.

Щоб мета цієї книги здійснилася, неодмінне тренування постави добровільно відпускати об'єкти Розуму. Медитувати не конче означає спорожнювати Розум, хоч це й можливий досвід. Під час медитації ми спостерігаємо появу об'єктів Розуму й дозволяємо їм іти, як буває з усіма явищами природи: вони постають і згасають.

Можна помітити, як сильно ми опираємося відпусканню певних думок, але якщо ми докладемо зусиль і справді виконаємо це завдання численні рази, користь буде колосальна. «Але ж це моє життя! Це мої проблеми!» Це саме той тип опору, який нам треба відпустити, бо можна жити в спокої без потреби ототожнюватися з думками. Я можу розв'язувати проблеми, але проблеми --- не я. Деякі люди розповідають, що під час медитації стають неспокійнішими, але правда в тому, що ці люди лише усвідомлюють, наскільки неспокійними вони вже є.

\begin{prat}{5, 10, 15, 20, 25 або 30 хвилин на день}
-- Відміряйте час за допомогою якогось звукового сигналу. Є чимало мобільних застосунків, здатних виконати це завдання;

-- Сядьте з прямим хребтом на сідничних горбах і покладіть руки на ноги;

-- Тримайте очі злегка розплющеними, а фокус розслабленим;

-- Не втручаючись, спостерігайте дихання;

-- Помітивши присутність будь-якої думки, лагідно поверніть увагу до дихання.
\end{prat}

\section*{Присутність}\addcontentsline{toc}{section}{Присутність}

У міру того, як практика медитації стає нам звичнішою, ми можемо поширити її на решту дня. Починаємо з простих справ, як-от ходьба, підмітання підлоги чи миття посуду. Тоді просуваємося до складніших справ, як-от робота, аж поки зможемо практикувати, спілкуючись з іншими людьми, під час розмов очі в очі або й в інтимні миті.

Фрустрація багатьох людей --- це спроба сприймати людський контакт як нагоду вразити іншого, і тому вони виставляють свої здобутки й мрії в пошуку схвалення. Але найвище, що ми можемо запропонувати, --- це наша цілковита присутність і увага.

Одна з великих красот стану присутності --- його нейтральність щодо станів розуму. Той, хто практикує старанно, не застрягає в негативних думках і не виснажується, шукаючи позитивних. Він лише розслабляється в рівновазі.

Стан Присутності також є розв'язком питання терпіння. Є два типи нетерплячих людей: ті, що визнають себе нетерплячими, і ті, що здаються терплячими. Але як можна страждати від нетерплячості, цілком поринувши в теперішню мить? Терплячий --- не той, хто терпить неспокій нетерплячості, а той, хто її не відчуває.

\begin{prat}{якнайчастіше}
-- Виконуйте справи, дихаючи усвідомлено й з повною увагою до акту їх виконання;

-- Якщо немає справи, яку слід зробити, виконуйте нічого-не-роблення з повною увагою.
\end{prat}

\section*{Задоволеність}\addcontentsline{toc}{section}{Задоволеність}

Поширена зумовленість --- відкладати щастя, плутаючи його з якимось конкретним досягненням чи задоволенням. І коли ми нарешті отримуємо бажане, відчуваємо велику ненаситну порожнечу й розуміємо, що досі схильні до того самого взірця.

Деякі люди мають «щастя» набути хвороби або низки хвороб, достатньо серйозних, щоб змусити їх переосмислити власні цінності й запросити приділити належну увагу багатствам, присутнім у кожній миті Життя. Але також на наше щастя нам не конче тяжко хворіти. Достатньо мати мужність усвідомити ці цикли незадоволення, які серйозно виснажили нашу планету.

Бути щасливим цієї миті може здатися контрінтуїтивним на перший погляд, але річ у тім, що неможливо бути щасливим в іншу мить. Справді, ця містична мить --- усе, що ми маємо.

\begin{prat}{якнайчастіше}
-- Хай навіть на кілька секунд, дайте собі право не потребувати думати, говорити чи рухатися й лише споглядайте те, що навколо вас.
\end{prat}

\chapter*{Про трансцендентну досконалість}\addcontentsline{toc}{chapter}{Про трансцендентну досконалість}

У дитинстві я часто впадав в екзистенційні кризи, бо не міг збагнути природу того, що звав «я». Від дуже раннього віку ця, здавалося б, невиправна туга гнітила мене, і я почувався так, ніби входжу у власну смерть. Чого я не знав, то це що я йшов у правильному напрямку.

Одразу зауважу, що неможливо говорити \textit{Істину}. Ми можемо говорити лише \textit{про} Істину. І роль цього розділу --- вказати на Неї.

Ще один важливий момент: суто інтелектуальне розуміння Істини не є насправді останнім знанням Істини. Істина --- не лише об'єкт Розуму.

Тож почнімо з уявлення, що нирка навчилася думати й каже «я» сама собі. Нирка може почати обурюватися тим, що серце безперервно переганяє їй кров для фільтрування. Нирка може також знайти причини бути щасливою, вдячною за кожен прожитий день. Але незалежно від того, що нирка думає про себе й про серце, вона є частиною більшого організму.

Чи ви колись думали, що було б, якби ліва рука почала почуватися гіршою за праву або права почала пишатися собою?

Тож можемо спитати, хто ми є, з більшою точністю.

\begin{prat}{раз на день}

-- Спитайте себе й споглядайте:

$\bullet$ Що є \textit{це}, що відчуває крізь цю шкіру?

$\bullet$ Що є \textit{це}, що спостерігає об'єкти розуму?

$\bullet$ Що є \textit{це}, що бачить цими очима й читає \textit{ці слова} саме зараз?
\end{prat}

Убого перекладаючи Тишу, \textit{це} і є весь Всесвіт, що бачить кожним оком і чує кожним вухом. Коли «я» чую «щось», Всесвіт чує Сам Себе. Коли «я» нюхаю «щось», Всесвіт нюхає Сам Себе. Саме цієї миті Всесвіт читає самого себе! Це максимум Істини, який я можу сказати. Якщо від мене вимагатимуть більшого, я сяду в медитацію й помру на якусь мить.

Той, хто має інтелектуальне розуміння Істини, покликаний до Праці. Той, хто мав містичний досвід Істини, лише втаємничений у Працю. Праця в її повному виконанні --- це втілена Істина: Єднання між Абсолютом та індивідом.

Духовно здійснені істоти не міркують категоріями правильного й неправильного, щоб керувати своїм життям, бо вони відчувають Єдність у своїх нутрощах. Для них більше не існує ні «мого щастя», ні «твого страждання». Ось чому ця книга не розглядає мораль і праведність. Я зробив так тому, що в моєму власному досвіді й розумінні праведне життя є тривіальним наслідком зосередженості. Той, хто розвивав старанне самоспостереження, відчуває в нутрощах незручність, що походить від брехні, ненависті, самозвеличання тощо, аж поки зрештою не вирішить із власної волі покинути такі звичаї.

У міру практикування запропонованих тут вправ ми відкриваємо благодать жити лагідно. Ми починаємо розуміти цінність найтонших переживань і дедалі менше потребуємо грубих. Випити чашку чаю може бути містичним переживанням!

Щасливий той, хто здатен увійти в Мир і дозволити Йому діяти у своєму житті.

\end{document}
